% SPDX-License-Identifier: CC-BY-NC-ND-4.0
% Copyright (c) 2026 Aymix — workbook content. See LICENSE-CONTENT.
% ============================ FRONT MATTER ============================
\newcommand{\frontheading}[2]{%
  \par\addvspace{7mm}%
  \noindent\begin{tikzpicture}[baseline=(t.base)]
    \node[fill=cBrand, text=white, rounded corners=1.3mm, font=\normalsize,
          inner sep=0pt, minimum width=8mm, minimum height=8mm] (n) {#1};
    \node[anchor=west, font=\sffamily\bfseries\LARGE, text=cInk] (t) at (n.east) [xshift=3mm] {#2};
  \end{tikzpicture}\par\vspace{1mm}%
  \noindent{\color{cBrand}\rule{\linewidth}{1pt}}\par\addvspace{2.5mm}%
}

% ---------------------------------------------------------------- COVER
\begin{titlepage}
\thispagestyle{empty}
\begin{tikzpicture}[remember picture, overlay]
  \fill[cBrandDeep] (current page.south west) rectangle (current page.north east);
  \fill[cBrand] ([yshift=-4.2cm]current page.north west) rectangle ([yshift=-9.7cm]current page.north east);
  \fill[cAccent] ([yshift=-4.2cm]current page.north west) rectangle ([yshift=-4.05cm]current page.north east);
  \fill[cAccent] ([yshift=-9.85cm]current page.north west) rectangle ([yshift=-9.7cm]current page.north east);
  % faint infrastructure glyph, lower-right
  \node[anchor=south east, text=white, opacity=0.04, font=\bfseries, scale=20]
     at ([xshift=0.7cm,yshift=-0.6cm]current page.south east) {\faCloud};
  % eyebrow
  \node[anchor=north west, text=cAccent, font=\sffamily\bfseries]
     at ([xshift=2.2cm,yshift=-2.5cm]current page.north west)
     {\Large BACKEND ENGINEERING\, \textbullet\, SPEEDRUN SERIES\, \textbullet\, VOLUME II};
  % title
  \node[anchor=north west, text=white, font=\sffamily\bfseries, align=left]
     at ([xshift=2.2cm,yshift=-5.25cm]current page.north west)
     {{\fontsize{39}{43}\selectfont The DevOps Engineer's}\\[3mm]
      {\fontsize{39}{43}\selectfont Lab Notebook}};
  % subtitle
  \node[anchor=north west, text=white, font=\sffamily, align=left]
     at ([xshift=2.2cm,yshift=-10.7cm]current page.north west)
     {\large A guided-practice companion to the};
  \node[anchor=north west, text=cAccent, font=\sffamily\bfseries, align=left]
     at ([xshift=2.2cm,yshift=-11.28cm]current page.north west)
     {\large DevOps Engineer --- 14-Day Speedrun Roadmap};
  % tagline
  \node[anchor=north west, text=white!85, font=\sffamily, align=left, text width=15.6cm]
     at ([xshift=2.2cm,yshift=-12.9cm]current page.north west)
     {From ``I can write a Dockerfile'' to designing, automating, securing, observing and operating
      production cloud infrastructure --- understanding \emph{why} each layer exists, not just which command to run.};
  % stat row
  \node[anchor=north west, align=left] at ([xshift=2.2cm,yshift=-15.3cm]current page.north west)
     {\begin{tikzpicture}[x=3.7cm, node distance=0pt]
        \foreach \i/\v/\l in {0/14/DAYS, 1/42/GUIDED LABS, 2/115+/HOURS DEPTH, 3/1/CAPSTONE}{
          \node[anchor=west, text=cAccent, font=\sffamily\bfseries] at (\i,0.35) {\Huge \v};
          \node[anchor=west, text=white!80, font=\sffamily\bfseries\footnotesize] at (\i,-0.35) {\l};
        }
      \end{tikzpicture}};
  % owner lines
  \node[anchor=north west, text=white!75, font=\sffamily, align=left]
     at ([xshift=2.2cm,yshift=-18.9cm]current page.north west)
     {This notebook belongs to:\ \ \textcolor{white!35}{\rule{7cm}{0.4pt}}};
  \node[anchor=north west, text=white!75, font=\sffamily, align=left]
     at ([xshift=2.2cm,yshift=-19.8cm]current page.north west)
     {Started on:\ \ \textcolor{white!35}{\rule{3.6cm}{0.4pt}}\qquad Target finish:\ \ \textcolor{white!35}{\rule{3.6cm}{0.4pt}}};
  \node[anchor=south west, text=white!55, font=\sffamily\footnotesize]
     at ([xshift=2.2cm,yshift=1.4cm]current page.south west)
     {Study the roadmap first \textbullet\ then learn with this workbook open beside you \textbullet\ break things on purpose};
\end{tikzpicture}
\end{titlepage}

% ---------------------------------------------------------------- HOW TO USE
\frontheading{\faCompass}{How to Use This Workbook}

This is not an exam-prep book and not a second reference manual. It is a \keyterm{laboratory manual} --- a
place to \emph{do the engineering} beside the roadmap. The roadmap tells you \emph{what} to learn; this
workbook walks you through understanding \emph{why it exists}, \emph{how it works underneath}, and
\emph{what it costs you} --- then makes you build, break and fix it.

The goal is not command recall. It is to leave each chapter reasoning like a senior engineer: able to
justify a design decision, name its trade-offs, and say what you would do differently at ten times the scale.

\wbsub{The daily study loop}
\begin{wbfigure}{Fig 0.1 — Run this loop once per day. The learning is in steps 3 and 4, not step 1.}
\wbscale{\begin{tikzpicture}[node distance=7mm and 12mm, every node/.style={font=\sffamily\footnotesize}]
  \node[wbbox, text width=30mm] (a) {\textbf{1 · Read}\\\scriptsize the roadmap day, first};
  \node[wbboxaccent, right=of a, text width=30mm] (b) {\textbf{2 · Work the 11 parts}\\\scriptsize concepts, diagrams, notes};
  \node[wbboxgood, right=of b, text width=30mm] (c) {\textbf{3 · Build it by hand}\\\scriptsize labs \& CloudBase};
  \node[wbboxwarn, below=8mm of c, text width=30mm] (d) {\textbf{4 · Break \& debug it}\\\scriptsize debugging labs};
  \node[wbboxops, below=8mm of a, text width=30mm] (e) {\textbf{5 · Reflect \& check}\\\scriptsize write-in spaces, checklist};
  \draw[wbarrow] (a)--(b); \draw[wbarrow] (b)--(c); \draw[wbarrow] (c)--(d);
  \draw[wbarrow] (d) -- (e); \draw[wbarrow] (e)--(a);
\end{tikzpicture}}
\end{wbfigure}

\wbsub{The eleven parts of every chapter}
\begin{center}\small
\begin{tabularx}{\linewidth}{@{}p{4mm}L{38mm}X@{}}
\textbf{1} & \textbf{The Big Picture} & What problem it solves, where it sits in a cloud architecture, how it links to earlier days.\\[1.2mm]
\textbf{2} & \textbf{Concept Map} & One page holding the whole territory and its vocabulary.\\[1.2mm]
\textbf{3} & \textbf{Guided Learning} & Each idea taught progressively: what / why it exists / how it works internally / when / when not / production example.\\[1.2mm]
\textbf{4} & \textbf{Visualizations} & Redraw the architecture from memory to lock it in.\\[1.2mm]
\textbf{5} & \textbf{Guided Hands-On Lab} & Build it step by step --- the workbook guides the decisions, you write the code.\\[1.2mm]
\textbf{6} & \textbf{Thinking Exercises} & Engineering reflections, not quizzes.\\[1.2mm]
\textbf{7} & \textbf{Debugging Practice} & A real incident worked as Observe → Hypothesise → Investigate → Root cause → Fix → Prevent.\\[1.2mm]
\textbf{8} & \textbf{Mini-Labs} & Three labs: Beginner, Intermediate, Production-style.\\[1.2mm]
\textbf{9} & \textbf{Engineering Notes} & Senior judgment, trade-offs, cost and reliability thinking.\\[1.2mm]
\textbf{10} & \textbf{Build-Along Project} & Extend \keyterm{CloudBase}, one cloud platform built across all 14 days.\\[1.2mm]
\textbf{11} & \textbf{Chapter Summary} & Takeaways, cheat sheet, CLI reference, official docs, production checklist.\\
\end{tabularx}
\end{center}

\begin{prodtip}
The write-in spaces and "redraw from memory" boxes are meant to be filled in \emph{by hand}. Print the pages
you are working, or annotate the PDF. Reproducing an architecture diagram yourself is what turns
"I have seen that" into "I can whiteboard that in an interview."
\end{prodtip}

\begin{secwarn}[on copy-paste]
This workbook deliberately does not hand you finished infrastructure. The labs guide you to build it. If you
paste a Terraform module or a Kubernetes manifest you did not type and reason through, you will feel
productive and learn nothing --- and you will not be able to debug it at 3am. The struggle in Parts 5, 7 and 10
is the entire point.
\end{secwarn}

% ---------------------------------------------------------------- LEGEND
\frontheading{\faInfoCircle}{Field Guide to the Engineering Notes}

Coloured callouts flag different \emph{kinds} of engineering thinking. Learn to read them at a glance.

\begin{multicols}{1}
\begin{seniornote}How an experienced engineer reasons about the decision --- the judgment behind the code.\end{seniornote}
\begin{prodtip}A concrete practice that holds up under real production load.\end{prodtip}
\begin{infranote}Something the infrastructure is doing underneath the abstraction.\end{infranote}
\begin{behind}What the tool or kernel is actually doing behind the command.\end{behind}
\begin{tradeoff}An honest design trade-off --- what you gain and what you give up.\end{tradeoff}
\begin{scalenote}What changes when traffic, nodes or teams multiply by ten.\end{scalenote}
\begin{reliability}A principle that keeps the system up when something fails.\end{reliability}
\begin{secwarn}A security trap --- read these twice.\end{secwarn}
\begin{costnote}Where the cloud bill quietly comes from, and how to avoid it.\end{costnote}
\begin{autotip}Where to replace human effort with automation.\end{autotip}
\begin{opsnote}Operational excellence: runbooks, on-call, and what "done" means.\end{opsnote}
\begin{drtip}Disaster recovery: what happens when a whole component or region is gone.\end{drtip}
\begin{perfnote}A latency or throughput consideration worth measuring.\end{perfnote}
\begin{interview}A question you will likely be asked, and how to answer it well.\end{interview}
\begin{misconception}A belief many engineers hold that is subtly wrong.\end{misconception}
\end{multicols}

\smallskip And the working blocks you will act on:

\begin{minilab}{Beginner / Intermediate / Production}{graded practice}
Labs with a Goal, Business context, Architecture, Requirements, Constraints, Hints, Expected outcome,
Production considerations and Common pitfalls.
\end{minilab}
\begin{debuglab}{A real production incident}
Symptom and evidence --- logs, \code{kubectl} output, a plan diff --- then a guided
Observe → Hypothesise → Investigate → Root cause → Fix → Prevent walkthrough.
\end{debuglab}
\begin{thinkbox}
An engineering reflection with ruled space to write your answer by hand.
\end{thinkbox}
\begin{buildalong}[CloudBase]
The one cloud platform you extend a little further every single day.
\end{buildalong}

% ---------------------------------------------------------------- THE ARC
\frontheading{\faMapSigns}{The 14-Day Arc at a Glance}

Each day earns its place and feeds the next. Read down the "CloudBase milestone" column to watch one
platform grow from a bare Linux host into production infrastructure with GitOps delivery and real SLOs.

\begin{center}\footnotesize
\wbrowcolors
\begin{tabularx}{\linewidth}{@{}C{7mm}L{38mm}X@{}}
\rowcolor{cBrand}\textcolor{white}{\textbf{Day}} & \textcolor{white}{\textbf{Topic}} & \textcolor{white}{\textbf{CloudBase milestone you build}}\\
1 & Linux \& the Command Line & Base host: strict-mode bootstrap, service user, systemd unit.\\
2 & Git \& Version Control & Repo conventions: branch protection, CODEOWNERS, secret-blocking hook.\\
3 & Networking for DevOps & Network design: /16 VPC, public/private subnets, bastion, least-privilege SGs.\\
4 & Docker Deep Dive & Hardened multi-stage image, non-root, minimal build context.\\
5 & Kubernetes Fundamentals & Deployment, Service, config/secrets, liveness \& readiness probes.\\
6 & Kubernetes in Production & Helm chart, Ingress with TLS, HPA, least-privilege RBAC.\\
7 & Terraform (IaC) & The network in code, with remote state and locking.\\
8 & Ansible (Config Mgmt) & Idempotent role replacing the bootstrap script; Vault secrets.\\
9 & CI/CD Engineering & test → scan → build → push, SHA-tagged, promoted not rebuilt.\\
10 & Cloud Infra on AWS & ALB, cluster, RDS in private subnets, least-privilege IAM.\\
11 & Monitoring \& Observability & Prometheus + Grafana, correlated logs, one real SLO and alert.\\
12 & Security \& Secrets & Secrets manager, image-scan gate, NetworkPolicies, RBAC audit.\\
13 & GitOps \& Deployment & Argo CD reconciliation, canary gated on SLO, rehearsed rollback.\\
14 & Capstone & Full assembly, DR drill, runbook, architecture review.\\
\end{tabularx}
\end{center}

\begin{keyidea}
\textbf{CloudBase} is deliberately a \emph{small} application carrying a \emph{large} amount of operational
weight --- exactly the shape of a strong take-home assignment or a well-scoped first project at a new job.
If you have worked Volume I of this series, the natural application to wrap in this infrastructure is the
order API you built there. Depth over breadth, every day.
\end{keyidea}

\clearpage
% ---------------------------------------------------------------- TOC
\frontheading{\faListUl}{Contents}
\vspace{1mm}
{\hypersetup{linkcolor=cInk}%
 \makeatletter\@starttoc{toc}\makeatother}
\vfill
\begin{center}\footnotesize\itshape\color{cInk!60}
Fourteen chapters \textbullet\ one growing platform \textbullet\ every layer understood, not just invoked.
\end{center}
\clearpage
